\documentclass[11pt]{article}
\usepackage[margin=1in]{geometry}
\usepackage{amsmath,amssymb,amsthm}
\usepackage{microtype}
\DeclareMathOperator{\tr}{tr}

\title{Error correction, superposition, and the $n=2$ response: a counting argument}
\date{}

\begin{document}
\maketitle

\noindent\textbf{What error correction requires.} If the residual stream is protected by an error-correcting scheme, the scheme has to distinguish feature directions --- along which perturbations are meaningful and should change the output --- from non-feature directions --- along which perturbations are noise and should be absorbed. Concretely: there should be a sensitivity threshold $\tau$ such that perturbations along features cross it (output changes) while perturbations along non-features stay below it (output does not change). Under superposition with $M > D$ features spanning $\mathbb{R}^D$, an arbitrary direction projects onto many features and a feature direction itself projects onto many others, so the discrimination cannot be done by looking at coefficients in an orthogonal decomposition --- it has to be done by a genuine sensitivity function $v \mapsto S(v)$ on the unit sphere.

\medskip
\noindent\textbf{What $n=2$ says about the sensitivity function.} Our empirical finding is that, to good approximation,
\[
S(v) \;\propto\; v^\top A v, \qquad A \in \mathbb{R}^{D\times D} \text{ symmetric.}
\]
A symmetric $A$ is determined by $D$ eigenvalues $\lambda_1 \geq \cdots \geq \lambda_D$ and a choice of orthonormal eigenbasis. The function $v \mapsto v^\top A v$ on $S^{D-1}$ is therefore fully described by these $D$ numbers (together with the eigenbasis). In particular, the set of directions above any threshold $\tau$,
\[
\mathcal{S}_\tau \;=\; \{\, v \in S^{D-1} : v^\top A v > \tau \,\},
\]
is the intersection of $S^{D-1}$ with an ellipsoid-complement: a geometrically simple, connected region whose shape is controlled by the eigenstructure of $A$.

\medskip
\noindent\textbf{The counting problem.} Error correction asks $\mathcal{S}_\tau$ to contain all $M$ features and as few non-features as possible. But:

\begin{itemize}
\item $\mathcal{S}_\tau$ is an open region of the sphere, not a discrete set; if it contains the features it contains their neighbourhoods too, so some non-features are inevitably admitted.
\item More substantively, the \emph{shape} of $\mathcal{S}_\tau$ has only $D$ degrees of freedom (the eigenvalues of $A$). Under $M > D$ features spanning $\mathbb{R}^D$, the features are spread across the sphere in a genuinely $D$-dimensional way. Asking a single ellipsoidal cap to simultaneously include all of them while excluding non-features over-constrains the $D$-dimensional freedom of $A$.
\end{itemize}

The sharpest form: if we demand every feature is \emph{maximally} sensitive ($d_j^\top A d_j = \lambda_{\max}(A)$ for every $j$), then all features lie in the top eigenspace of $A$. Under features spanning $\mathbb{R}^D$, the top eigenspace must equal $\mathbb{R}^D$, which forces $A \propto I$. But $A \propto I$ means $v^\top A v$ is constant on the sphere --- every direction is equally sensitive, so \emph{no} direction is privileged and there is nothing left for error correction to do.

\medskip
\noindent\textbf{Contrast with $n \to \infty$.} This counting obstruction is specific to the quadratic case. At $n \to \infty$ the sensitivity is $S(v) = \max_i |d_i^\top v|/\tau_i$, a piecewise-linear function of $v$ whose level sets \emph{are} unions of feature-centred caps, one per feature. A piecewise-linear max-of-$M$ function has $M$ local maxima by construction, so $M > D$ features can all be individually privileged. That is exactly why per-feature threshold error correction lives at $n = \infty$ and predicts square iso-plateaus. The observation $n \approx 2$ is therefore not merely inconsistent with one particular EC scheme: the functional form has the wrong number of degrees of freedom to support any EC that privileges $M > D$ distinct features.

\medskip
\noindent\textbf{Summary.} Under superposition with features spanning $\mathbb{R}^D$, error correction needs a sensitivity function that can single out $M > D$ directions from the rest of the sphere. A quadratic form on $S^{D-1}$ cannot do this: it has only $D$ eigen-parameters, and the region above any threshold is a single ellipsoidal cap. Our empirical $n \approx 2$ forces sensitivity to be (approximately) such a quadratic form, so the model's activation stream is not equipped with any sensitivity structure that discriminates $M > D$ feature directions from their ambient non-feature directions. The conclusion is stronger than ``per-feature threshold EC is ruled out'': any error correction that acts by elevating feature directions above non-feature directions, at the level of response sensitivity, is incompatible with the observed $n = 2$ under superposition.

\end{document}
